\documentclass[11pt]{amsart}
\usepackage[T1]{fontenc}
\usepackage[margin=1.02in]{geometry}
\DeclareFontShape{OT1}{cmr}{bx}{n}{<-> cmbx10}{}
\DeclareMathSizes{8}{8}{8}{8}
\DeclareMathSizes{9}{9}{8}{8}
\DeclareMathSizes{10.95}{10.95}{8}{8}
\newtheorem{theorem}{Theorem}
\newcommand{\kl}{\operatorname{kl}}
\newcommand{\KL}{\operatorname{KL}}
\newcommand{\Bin}{\operatorname{Bin}}
\newcommand{\E}{\operatorname{\mathbb E}}
\title{A strict gap for a common PAC-Bayes comparator}
\date{November 27, 2026}
\begin{document}
\begin{abstract}
The pointwise optimistic PAC-Bayes-kl inversion need not be
attainable in expectation by one comparator. Two independent
Bernoulli observations give a strict expectation gap of at least
\((7-\sqrt{34})/80\) for every proper closed convex comparator in
the generic theorem. The obstruction is a shared binomial moment
budget across two possible samples. A fixed Catoni comparator
bounds the gap from above by \(9/176\) on this instance. This
answers the universal attainability part of Open Problem~1 of
Foong, Bruinsma, Burt and Turner.
\end{abstract}
\maketitle

\section{Setup and the pointwise floor}
For \(N\) independent observations and a bounded loss in \([0,1]\),
let \(q\) denote the posterior-averaged empirical risk. Write
\(K=\KL(Q\Vert P)\) for posterior-to-prior divergence and
\(a=(K+\log(1/\delta))/N\), with \(\delta\in(0,1)\).
Foong et al.\ \cite[Theorem~3 and Eq.~(6)]{foong} state the
generic PAC-Bayes inequality for a proper closed convex
\(\Delta:[0,1]^2\to\mathbb R\cup\{+\infty\}\), chosen before
sampling. Its moment factor and upper inversion are
\[
 I=\sup_{r\in[0,1]}\E_{B\sim\Bin(N,r)} e^{N\Delta(B/N,r)},
 \qquad b_\Delta=\sup\{p\in[0,1]:\Delta(q,p)\le a+\log I/N\},
\]
with the empty sublevel set giving 1. A moment is taken over the
support of the binomial law. Infinite \(I\) gives the vacuous
bound 1. The pointwise ideal uses \(b_*=\sup\{p:\kl(q,p)\le a\}\),
where \(\kl\) is Bernoulli relative entropy. Foong et al.\
\cite[Open Problem~1]{foong} ask whether some one \(\Delta\)
attains this floor in expectation for an arbitrary dataset law.

With finite \(I\), the positive grid masses at interior \(r\)
give \(I>0\). Normalize \(f=\Delta-\log I/N\), so
\[
            \E_{B\sim\Bin(N,r)} e^{Nf(B/N,r)}\le1.
\]
For a fixed interior \(p\), convexity in the first argument, followed
by the finite change-of-measure inequality, gives
\[
\begin{split}
 N f(q,p)&\le \E_{\Bin(N,q)} Nf(B/N,p)\\
          &\le\log\E_{\Bin(N,p)} e^{Nf(B/N,p)}
                   +\KL(\Bin(N,q)\Vert\Bin(N,p))
          \le N\kl(q,p).
\end{split}
\]
The second inequality follows by applying Jensen's inequality for
the logarithm under the tilted law after dividing by its density
relative to \(\Bin(N,p)\) on its support. It also holds when
\(q=0\) or \(q=1\). Thus every interior point of the ideal
sublevel set is admissible under \(f\), and approximation from
inside \((0,1)\) gives \(b_\Delta\ge b_*\), including endpoint
suprema. Indeed \(a>0\) ensures the ideal sublevel set contains
interior points and its supremum is approachable through them.
This recovers the pointwise comparison in \cite[Remark~2]{foong}.

\section{A two-sample obstruction}
\begin{theorem}
There is a setup with two iid observations, two hypotheses and
\(\delta=3/4\) for which
\[
  \frac{7-\sqrt{34}}{80}
     \le \inf_{\Delta}\E (b_\Delta-b_*) \le \frac9{176}.
\]
In particular the ideal is unattainable even in the infimum.
\end{theorem}

\begin{proof}
Take inputs supported at a single point and iid fair labels.
Use zero-one loss and two hypotheses predicting zero at that
point, with different predictions at another point of probability
zero. Put prior mass \(1/12\) on the first hypothesis. The posterior
puts all mass on it when both labels are zero, and equals the
prior otherwise. All four ordered label samples are equally likely,
and the induced values are
\[
\begin{array}{c|c|c|c}
  q&K&\text{probability}&b_*\\ \hline
  0&\log12&1/4&3/4\\
  1/2&0&1/2&3/4\\
  1&0&1/4&1
\end{array}
\]
To see the first two inversions, the respective thresholds are
\(a_0=\log4\) and \(a_1=\log(4/3)/2\), and at \(p=3/4\) the
exponentiated doubled divergences are \(16\) and \(4/3\).
The divergence is increasing on the upper branch \(p>q\).
The \(q=1\) row always has \(b_\Delta=1\) by the pointwise floor.

For finite \(I\), suppose that both generic inversions of the
first two rows are strictly below an interior \(r\).
Then \(f(0,r)>a_0\) and \(f(1/2,r)>a_1\), because otherwise \(r\)
would be admissible in the corresponding row. The first two
terms of the normalized binomial moment at \(r\) therefore have
sum strictly greater than
\[
      L(r)=16(1-r)^2+\frac{8}{3}r(1-r).
\]
At \(r=4/5\), \(L(r)=16/15>1\), contradicting the moment bound.
The same argument works at the smaller root of \(L(r)=1\),
which lies between \(4/5\) and \(1\). In fact,
\(3(L(r)-1)=40r^2-88r+45\), giving \(r=(22-\sqrt{34})/20\).
Thus one of the two inversions is at least \(r\), and its row has
probability at least \(1/4\). The pointwise floor precludes
cancellation by the other rows. The gap is at least
\((r-3/4)/4=(7-\sqrt{34})/80>1/80\).
The same lower bound holds with the vacuous inversion for
infinite \(I\).

For the upper bound, take the Catoni comparator
\[
 C_\beta(q,p)=-\beta q-\log(1-p+p e^{-\beta}),\qquad \beta=\log12 .
\]
It is closed convex and its binomial moment equals 1 by the
binomial generating function \cite[Corollary~1]{foong}.
In the first two rows the inversion condition simplifies to
\(1-11p/12\ge1/4\), giving \(9/11\) in both.
The third row gives 1. The resulting expected excess is
\((3/4)(9/11-3/4)=9/176\).
\end{proof}

\section{A distribution-free comparison}
A direct upper bound follows from the ordinary \(\Delta=\kl\)
choice. For interior \(r\), the binomial moment \(I_{\kl}\)
is the sum over \(k=0,\ldots,N\) of the mass at \(k\) in
\(\Bin(N,k/N)\), by cancelling the \(r\)-dependent factors.
At the endpoints the moment is 1. Hence \(I_{\kl}\le N+1\).
Sharper moment estimates are classical \cite[Corollary~2]{foong}.

If \(b_*=1\), the excess is zero. Otherwise \(b_*\ge q\) and
\(\kl(q,b_*)=a\). For interior \(p\ge b_*\), the derivative with
respect to \(p\) satisfies
\[
      \frac{p-q}{p(1-p)}\ge 4(p-b_*).
\]
Integrating gives \(\kl(q,p)-\kl(q,b_*)\ge2(p-b_*)^2\), also
at a boundary by limits. The ordinary KL inversion increases the
level \(a\) by \(\log I_{\kl}/N\). Its excess over \(b_*\) is at most
\(\sqrt{\log I_{\kl}/(2N)}\), uniformly in the data. In particular
the best common comparator's expectation gap is at most
\(\sqrt{\log(N+1)/(2N)}\) in every setup. The explicit example
above does not determine the exact infimum.

\section*{Scope, verification and technical disclosure}
Theorem~4 of \cite{foong} optimizes for fixed \(q,K\), while
Corollary~3 supplies a non-strict expectation floor. The
numerical comparisons in their Figure~2(d) address the Catoni
family. Hellstr\"om and Guedj \cite[Theorem~8]{hg} distinguish
the moment-one benchmark from the valid bound retaining its
moment factor. The obstruction here concerns the expectation
of samplewise inversions for one comparator in the generic
theorem. The validity of the optimistic high-probability bound
by other means is a separate question.

The Python replay checks the finite law and the exact moment
arithmetic. A partial Lean~4 certificate checks the polynomial
identity and rational threshold algebra. The analytical proof,
including Jensen's inequality and the irrational root step,
is written here. Internally reviewed, not peer reviewed.

Under human direction, ChatGPT Astra~6.2 (OpenAI) generated the
discovery proof and research packet. OpenAI Codex independently
checked the mathematics. GPT-6.2 Sol (OpenAI) checked the
literature and prepared the manuscript, partial Lean certificate
and illustrations. DannyExperiments maintains the repository
and is the preferred citation identity. No AI system is proposed
as an author.

\begin{thebibliography}{9}
\raggedright
\bibitem{foong} A. Y. K. Foong, W. P. Bruinsma, D. R. Burt and R. E. Turner,
\emph{How tight can PAC-Bayes be in the small data regime?},
Advances in Neural Information Processing Systems 34 (2021).
\bibitem{hg} F. Hellstr\"om and B. Guedj,
\emph{Comparing comparators in generalization bounds},
Proc. AISTATS, PMLR 238 (2024), 73--81.
\end{thebibliography}
\end{document}
